\subsection{超滤子}

定义 4. 集合 X 上的超滤子是这样的滤子 F：X 上不存在严格细于 F 的滤子（换言之，即 X 上所有滤子之有序集中的极大元素）。

由于 X 上所有滤子之有序集是归纳的（第 2 目，命题 1，推论 3），Zorn 引理（《集合论》，R，§ 6，第 10 目）表明：

定理 1. 若 F 是集合 X 上的任一滤子，则存在细于 F 的超滤子。

命题 5. 设 \(\mathfrak{F}\) 是集合 \(\mathbf{X}\) 上的超滤子。若 \(\mathbf{A}\) 与 \(\mathbf{B}\) 是 \(\mathbf{X}\) 的两个子集，且 \(\mathbf{A} \cup \mathbf{B} \in \mathfrak{F}\)，则或者 \(\mathbf{A} \in \mathfrak{F}\)，或者 \(\mathbf{B} \in \mathfrak{F}\)。

若命题不真，则存在 X 的子集 A 与 B，使得 A∉F、B∉F 且 A∪B∈F。设 G 是 X 中满足 A∪M∈F 的子集 M 所成之集。直接验证可知 G 是 X 上的滤子，且 G 严格细于 F，因为 B∈G；但这与 F 是超滤子的假设矛盾。

推论. 若子集的一个有限序列 \((\mathbf{A}_i)_{1 \leqslant i \leqslant n}\)（\(\mathbf{X}\) 的）之并属于超滤子 \(\mathfrak{F}\)，则至少有一个 \(\mathbf{A}_i\) 属于 \(\mathfrak{F}\)。

对 n 用归纳法证明。

特别地，若 \((\mathbf{A}_i)_{1\leqslant i\leqslant n}\) 是 \(\mathbf{X}\) 的一个覆盖，则至少有一个 \(\mathbf{A}_i\) 属于 \(\mathfrak{F}\)。

命题 5 刻画了超滤子；更一般地，我们有：

命题 6. 设 \(\mathfrak{G}\) 是集合 X 上某个滤子的子基。若对 X 的每个子集 Y，或者 \(Y \in \mathfrak{G}\)，或者 \(\complement Y \in \mathfrak{G}\)，则 \(\mathfrak{G}\) 是 X 上的超滤子。

设 \(\mathfrak{F}\) 是包含 \(\mathfrak{G}\) 的滤子（按假设它是存在的）；则 \(\mathfrak{F}\) 与 \(\mathfrak{G}\) 一致；因为若 \(Y \in \mathfrak{F}\)，则 \(\complement Y \notin \mathfrak{F}\)；于是 \(\complement Y \notin \mathfrak{G}\)，从而 \(Y \in \mathfrak{G}\)。

超滤子的例子. 非空集合 X 中包含给定元素 \(a \in X\) 的所有子集所成之集是一个超滤子；因为它是滤子，且若 Y 是 X 的任一子集，则或者 \(a \in Y\)，或者 \(a \in \complement Y\)。这样的超滤子称为平凡超滤子。

除这个例子之外，我们将只在用到定理 1（从而用到选择公理）时才证明超滤子（即使在可数无限集上的超滤子）的存在性。

注. 若 X 至少含两个元素，则 X 上至少有两个不同的超滤子，因而 X 上滤子的有序集没有最大元。

命题 7. 每个滤子 \(\mathfrak{F}\)（集合 \(X\) 上的）都是细于 \(\mathfrak{F}\) 的诸超滤子之交。

显然这个交包含 \(\mathfrak{F}\)。反之，设 \(A\) 是 \(X\) 的不属于 \(\mathfrak{F}\) 的子集，并用 \(A'\) 表示 \(\complement A\)；\(A\) 不包含 \(\mathfrak{F}\) 的任何集合；于是每个 \(M \in \mathfrak{F}\) 都与 \(A'\) 相交，从而（第 2 目，命题 1，推论 1）存在滤子 \(\mathfrak{F}'\)，它细于 \(\mathfrak{F}\) 且包含 \(A'\)。若 \(\mathfrak{U}\) 是细于 \(\mathfrak{F}'\) 的超滤子（定理 1），则由此可知 \(A \notin \mathfrak{U}\)。证明完毕。

